\chapter{2007年大纲I卷（文）}

\section{单选题}

\begin{problem}
设 \(S = \{x \mid 2x + 1 > 0\}\)，\(T = \{x \mid 3x - 5 < 0\}\)，则 \(S \cap T =\)（\quad）

\choices
    {\(\varnothing\)}
    {\(\left\{x\mid x < -\frac{1}{2}\right\}\)}
    {\(\left\{x\mid x > \frac{5}{3}\right\}\)}
    {\(\left\{x\mid -\frac{1}{2} < x < \frac{5}{3}\right\}\)}
\end{problem}

\begin{answer}
D
\end{answer}

\begin{solution}
由 \(2x+1>0\) 得 \(x>-\frac{1}{2}\)，由 \(3x-5<0\) 得 \(x<\frac{5}{3}\). 因此
\(S\cap T=\left\{x\mid -\frac{1}{2}<x<\frac{5}{3}\right\}\)，选择 D.
\end{solution}

\begin{problem}
\(\alpha\) 是第四象限角，\(\cos \alpha = \frac{12}{13}\)，则 \(\sin \alpha =\)（\quad）

\choices
    {\(\frac{5}{13}\)}
    {\(-\frac{5}{13}\)}
    {\(\frac{5}{12}\)}
    {\(-\frac{5}{12}\)}
\end{problem}

\begin{answer}
B
\end{answer}

\begin{solution}
因为 \(\alpha\) 是第四象限角，所以 \(\sin\alpha<0\). 由
\(\sin^2\alpha+\cos^2\alpha=1\)，得
\(\sin\alpha=-\sqrt{1-\left(\frac{12}{13}\right)^2}=-\frac{5}{13}\)，选择 B.
\end{solution}

\begin{problem}
已知向量 \(\bs{a} = (-5, 6)\)，\(\bs{b} = (6, 5)\)，则 \(\bs{a}\) 与 \(\bs{b}\)（\quad）

\choices
    {垂直}
    {不垂直也不平行}
    {平行且同向}
    {平行且反向}
\end{problem}

\begin{answer}
A
\end{answer}

\begin{solution}
\(\bs{a}\cdot\bs{b}=(-5)\times6+6\times5=0\). 由于两个向量均为非零向量，故 \(\bs{a}\perp\bs{b}\)，选择 A.
\end{solution}

\begin{problem}
已知双曲线的离心率为 \(2\)，焦点是 \((-4, 0)\)，\((4, 0)\)，则双曲线方程为（\quad）

\choices
    {\(\frac{x^2}{4} - \frac{y^2}{12} = 1\)}
    {\(\frac{x^2}{12} - \frac{y^2}{4} = 1\)}
    {\(\frac{x^2}{10} - \frac{y^2}{6} = 1\)}
    {\(\frac{x^2}{6} - \frac{y^2}{10} = 1\)}
\end{problem}

\begin{answer}
A
\end{answer}

\begin{solution}
设双曲线的实半轴长为 \(a\)，虚半轴长为 \(b\)，焦半距为 \(c\). 由焦点坐标得 \(c=4\)，由离心率得
\(\frac{c}{a}=2\)，所以 \(a=2\). 又 \(c^2=a^2+b^2\)，从而
\(b^2=4^2-2^2=12\). 因此双曲线方程为
\(\frac{x^2}{4}-\frac{y^2}{12}=1\)，选择 A.
\end{solution}

\begin{problem}
甲，乙，丙 \(3\text{ 位}\)同学选修课程. 从 \(4\text{ 门}\)课程中，甲选修 \(2\text{ 门}\)，乙，丙各选修 \(3\text{ 门}\)，则不同的选修方案共有（\quad）

\choices
    {\(36\text{ 种}\)}
    {\(48\text{ 种}\)}
    {\(96\text{ 种}\)}
    {\(192\text{ 种}\)}
\end{problem}

\begin{answer}
C
\end{answer}

\begin{solution}
甲的选法有 \(\mathrm C_4^2\) 种，乙、丙的选法各有 \(\mathrm C_4^3\) 种. 三人的选课相互独立，所以方案总数为
\(\mathrm C_4^2\mathrm C_4^3\mathrm C_4^3=6\times4\times4=96\)，选择 C.
\end{solution}

\begin{problem}
下面给出的四个点中，位于 \(\begin{cases}
x + y - 1 < 0,\\
x - y + 1 > 0
\end{cases}\) 表示的平面区域内的点是（\quad）

\choices
    {\((0,2)\)}
    {\((-2,0)\)}
    {\((0, -2)\)}
    {\((2,0)\)}
\end{problem}

\begin{answer}
C
\end{answer}

\begin{solution}
逐点检验两个不等式. 点 \((0,2)\) 不满足 \(x+y-1<0\)；点 \((-2,0)\) 的
\(x-y+1=-1\)，不满足第二个不等式；点 \((0,-2)\) 满足
\(0-2-1=-3<0\) 及 \(0-(-2)+1=3>0\)；点 \((2,0)\) 不满足第一个不等式. 因此只有点 \((0,-2)\) 在该区域内，选择 C.
\end{solution}

\begin{problem}
如图，正四棱柱 \(ABCD - A_{1}B_{1}C_{1}D_{1}\) 中，\(AA_{1} = 2AB\)，则异面直线 \(A_{1}B\) 与 \(AD_{1}\) 所成角的余弦值为（\quad）

\bitmapfigure[max width=0.42\linewidth]{img/2007/syllabus_paper_1_liberal/q7_fig1.png}

\choices
    {\(\frac{1}{5}\)}
    {\(\frac{2}{5}\)}
    {\(\frac{3}{5}\)}
    {\(\frac{4}{5}\)}
\end{problem}

\begin{answer}
D
\end{answer}

\begin{solution}
设正方形底面边长为 \(s\)，建立空间直角坐标系，使
\(A=(0,0,0)\)，\(B=(s,0,0)\)，\(D=(0,s,0)\)，\(A_1=(0,0,2s)\)，\(D_1=(0,s,2s)\). 于是
\(\overrightarrow{A_1B}=(s,0,-2s)\)，\(\overrightarrow{AD_1}=(0,s,2s)\).

两向量的内积为 \(-4s^2\)，且长度均为 \(s\sqrt{5}\). 异面直线所成的角取锐角，故其余弦为
\(\frac{\left|-4s^2\right|}{s\sqrt{5}\times s\sqrt{5}}=\frac{4}{5}\)，选择 D.
\end{solution}

\begin{problem}
设 \(a > 1\)，函数 \(f(x) = \log_a x\) 在区间 \([a, 2a]\) 上的最大值与最小值之差为 \(\frac{1}{2}\)，则 \(a =\)（\quad）

\choices
    {\(\sqrt{2}\)}
    {\(2\)}
    {\(2\sqrt{2}\)}
    {\(4\)}
\end{problem}

\begin{answer}
D
\end{answer}

\begin{solution}
因为 \(a>1\)，函数 \(f(x)=\log_a x\) 在正数范围内单调递增，所以在 \([a,2a]\) 上的最小值和最大值分别为
\(f(a)=1\) 和 \(f(2a)=1+\log_a2\). 由题意
\(\log_a2=\frac{1}{2}\)，即 \(a^{1/2}=2\)，所以 \(a=4\)，选择 D.
\end{solution}

\begin{problem}
\(f(x)\)，\(g(x)\) 是定义在 \(\R\) 上的函数，\(h(x) = f(x) + g(x)\)，则“\(f(x)\)，\(g(x)\) 均为偶函数”是“\(h(x)\) 为偶函数”的（\quad）

\choices
    {充要条件}
    {充分而不必要的条件}
    {必要而不充分的条件}
    {既不充分也不必要的条件}
\end{problem}

\begin{answer}
B
\end{answer}

\begin{solution}
若 \(f(x)\)，\(g(x)\) 均为偶函数，则
\(h(-x)=f(-x)+g(-x)=f(x)+g(x)=h(x)\)，所以这是 \(h(x)\) 为偶函数的充分条件.

但它不是必要条件. 例如取 \(f(x)=x\)，\(g(x)=-x\)，则 \(f\)，\(g\) 均不是偶函数，而 \(h(x)=0\) 是偶函数. 因此选择 B.
\end{solution}

\begin{problem}
函数 \( y = 2\cos^2 x \) 的一个单调增区间是（\quad）

\choices
    {\(\left(-\frac{\pi}{4}, \frac{\pi}{4}\right)\)}
    {\(\left(0, \frac{\pi}{2}\right)\)}
    {\(\left(\frac{\pi}{4}, \frac{3\pi}{4}\right)\)}
    {\(\left(\frac{\pi}{2}, \pi\right)\)}
\end{problem}

\begin{answer}
D
\end{answer}

\begin{solution}
\(y=2\cos^2x=1+\cos2x\)，所以 \(y'=-2\sin2x\). 在
\(\left(\frac{\pi}{2},\pi\right)\) 内有 \(\sin2x<0\)，从而 \(y'>0\)，函数在该区间上单调递增. 因此选择 D.
\end{solution}

\begin{problem}
曲线 \( y = \frac{1}{3}x^{3} + x \) 在点 \(\left(1, \frac{4}{3}\right)\) 处的切线与坐标轴围成的三角形面积为（\quad）

\choices
    {\(\frac{1}{9}\)}
    {\(\frac{2}{9}\)}
    {\(\frac{1}{3}\)}
    {\(\frac{2}{3}\)}
\end{problem}

\begin{answer}
A
\end{answer}

\begin{solution}
曲线的导数为 \(y'=x^2+1\)，故在 \(x=1\) 处的切线斜率为 \(2\). 切线方程为
\(y-\frac{4}{3}=2(x-1)\)，即 \(y=2x-\frac{2}{3}\).

令 \(y=0\)，得横轴截距为 \(\frac{1}{3}\)；令 \(x=0\)，得纵轴截距为 \(-\frac{2}{3}\). 所以三角形面积为
\(\frac{1}{2}\times\frac{1}{3}\times\frac{2}{3}=\frac{1}{9}\)，选择 A.
\end{solution}

\begin{problem}
抛物线 \( y^{2} = 4x \) 的焦点为 \(F\)，准线为 \(l\)，经过 \(F\) 且斜率为 \(\sqrt{3}\) 的直线与抛物线在 \(x\) 轴上方的部分相交于点 \(A\)，\(AK \perp l\)，垂足为 \(K\)，且 \(\triangle AKF\) 的面积是（\quad）

\choices
    {\(4\)}
    {\(3\sqrt{3}\)}
    {\(4\sqrt{3}\)}
    {\(8\)}
\end{problem}

\begin{answer}
C
\end{answer}

\begin{solution}
抛物线 \(y^2=4x\) 的焦点为 \(F=(1,0)\)，准线为 \(l:x=-1\). 过 \(F\) 的直线方程为
\(y=\sqrt{3}(x-1)\). 与抛物线联立，得
\(3(x-1)^2=4x\)，即 \(3x^2-10x+3=0\)，所以 \(x=\frac{1}{3}\) 或 \(x=3\). 取位于 \(x\) 轴上方的交点，得 \(A=(3,2\sqrt{3})\).

由于 \(AK\perp l\)，且 \(l\) 为竖直直线，所以 \(K=(-1,2\sqrt{3})\). 于是 \(AK=4\)，点 \(F\) 到直线 \(AK\) 的距离为 \(2\sqrt{3}\)，从而
\([\triangle AKF]=\frac{1}{2}\times4\times2\sqrt{3}=4\sqrt{3}\)，选择 C.
\end{solution}

\section{填空题}

\begin{problem}
从某自动包装机包装的食盐中，随机抽取 \(20\text{ 袋}\)，测得各袋的质量分别为（单位：\(\mathrm{g}\)）：

\(492\)，\(496\)，\(494\)，\(495\)，\(498\)，\(497\)，\(501\)，\(502\)，\(504\)，\(496\)，\(497\)，\(503\)，\(506\)，\(508\)，\(507\)，\(492\)，\(496\)，\(500\)，\(501\)，\(499\).

根据频率分布估计总体分布的原理，该自动包装机包装的袋装食盐质量在 \(497.5\mathrm{~g}\sim 501.5\mathrm{~g}\) 之间的概率约为\fillinblank{}.
\end{problem}

\begin{answer}
\(0.25\)
\end{answer}

\begin{solution}
在区间 \(497.5\mathrm{~g}\sim501.5\mathrm{~g}\) 内的整数质量为 \(498\mathrm{~g}\)，\(499\mathrm{~g}\)，\(500\mathrm{~g}\)，\(501\mathrm{~g}\). 样本中它们分别出现 \(1\)，\(1\)，\(1\)，\(2\) 次，共 \(5\) 袋. 因此用样本频率估计总体概率，所得概率约为
\(\frac{5}{20}=0.25\).
\end{solution}

\begin{problem}
函数 \( y = f(x) \) 的图象与函数 \( y = \log_{3} x \)（\(x > 0\)）的图象关于直线 \( y = x \) 对称，则 \( f(x) = \)\fillinblank{}.
\end{problem}

\begin{answer}
\(3^{x}\)（\(x\in\R\)）
\end{answer}

\begin{solution}
关于直线 \(y=x\) 对称的两个函数图象互为反函数. 由
\(y=\log_3x\) 得 \(x=3^y\)，交换 \(x\)，\(y\) 后得到反函数
\(y=3^x\)，其定义域为 \(\R\).
\end{solution}

\begin{problem}
正四棱锥 \(S - ABCD\) 的底面边长和各侧棱长都为 \(\sqrt{2}\)，点 \(S\)，\(A\)，\(B\)，\(C\)，\(D\) 都在同一个球面上，则该球的体积为\fillinblank{}.
\end{problem}

\begin{answer}
\(\frac{4\pi}{3}\)
\end{answer}

\begin{solution}
设正方形底面中心为 \(O\)，正四棱锥的高为 \(h\). 底面边长为 \(\sqrt{2}\)，所以底面中心到顶点的距离为
\(OA=\frac{\sqrt{2}}{\sqrt{2}}=1\). 又侧棱长为 \(\sqrt{2}\)，在直角三角形 \(SOA\) 中
\(h^2+OA^2=SA^2\)，即 \(h^2+1=2\)，故 \(h=1\).

底面四个顶点到 \(O\) 的距离均为 \(1\)，且 \(OS=h=1\)，所以以 \(O\) 为球心、半径为 \(1\) 的球经过 \(S\)，\(A\)，\(B\)，\(C\)，\(D\). 该球的体积为
\(\frac{4}{3}\pi\times1^3=\frac{4\pi}{3}\).
\end{solution}

\begin{problem}
等比数列 \(\{a_{n}\}\) 的前 \(n\) 项和 \(S_{n}\)，已知 \(S_{1}\)，\(2S_{2}\)，\(3S_{3}\) 成等差数列，则 \(\{a_{n}\}\) 的公比为\fillinblank{}.
\end{problem}

\begin{answer}
\(\frac{1}{3}\)
\end{answer}

\begin{solution}
设等比数列的首项为 \(a_1\)，公比为 \(q\). 由等比数列的定义，\(a_1\ne0\)，且 \(q\ne0\). 则
\(S_1=a_1\)，\(S_2=a_1(1+q)\)，\(S_3=a_1(1+q+q^2)\).
因为 \(S_1\)，\(2S_2\)，\(3S_3\) 成等差数列，所以
\(2\times2S_2=S_1+3S_3\). 代入上式并约去非零的 \(a_1\)，得
\(4(1+q)=1+3(1+q+q^2)\)，即 \(q(1-3q)=0\).

等比数列的公比 \(q\ne0\)，所以 \(q=\frac{1}{3}\).
\end{solution}

\section{解答题}

\begin{problem}
设锐角三角形 \(ABC\) 的内角 \(A\)，\(B\)，\(C\) 的对边分别为 \(a\)，\(b\)，\(c\)，\(a = 2b\sin A\).

\begin{enumerate}
\item 求 \(B\) 的大小；
\item 若 \(a = 3\sqrt{3}\)，\(c = 5\)，求 \(b\).
\end{enumerate}
\end{problem}

\begin{solution}
\begin{enumerate}
\item 由正弦定理，\(\frac{a}{\sin A}=\frac{b}{\sin B}\)，所以
\(a=\frac{b\sin A}{\sin B}\). 结合题设 \(a=2b\sin A\)，并注意到锐角三角形中 \(b>0\)，\(\sin A>0\)，可得
\(\sin B=\frac{1}{2}\). 又因为 \(B\) 为锐角，所以 \(B=30^{\circ}\).
\item 由余弦定理，\(b^2=a^2+c^2-2ac\cos B\). 代入
\(a=3\sqrt{3}\)，\(c=5\)，\(B=30^{\circ}\)，得
\(b^2=27+25-2\times3\sqrt{3}\times5\times\frac{\sqrt{3}}{2}=7\). 因为边长为正，所以 \(b=\sqrt{7}\).
\end{enumerate}
\end{solution}

\begin{problem}
某商场经销某商品，顾客可采用一次性付款或分期付款购买. 根据以往资料统计，顾客采用一次性付款的概率是 \(0.6\). 经销一件该商品，若顾客采用一次性付款，商场获得利润 \(200\text{ 元}\)；若顾客采用分期付款，商场获得利润 \(250\text{ 元}\).

\begin{enumerate}
\item 求 \(3\text{ 位}\)购买该商品的顾客中至少有 \(1\text{ 位}\)采用一次性付款的概率；
\item 求 \(3\text{ 位}\)顾客每人购买 \(1\text{ 件}\)该商品，商场获得利润不超过 \(650\text{ 元}\)的概率.
\end{enumerate}
\end{problem}

\begin{solution}
\begin{enumerate}
\item 每位顾客采用分期付款的概率为 \(1-0.6=0.4\). 在各位顾客的选择相互独立的条件下，三位顾客均采用分期付款的概率为 \(0.4^3\). 因此至少有一位采用一次性付款的概率为
\(1-0.4^3=0.936\).
\item 设三位顾客中采用一次性付款的人数为 \(X\)，则总利润为
\(200X+250(3-X)=750-50X\). 利润不超过 \(650\) 元等价于
\(750-50X\leqslant650\)，即 \(X\geqslant2\). 所以所求概率为
\(\mathrm C_3^2(0.6)^2(0.4)+(0.6)^3=0.648\).
\end{enumerate}
\end{solution}

\begin{problem}
四棱锥 \(S - ABCD\) 中，底面 \(ABCD\) 为平行四边形，侧面 \(SBC \perp\) 底面 \(ABCD\). 已知 \(\angle ABC = 45^{\circ}\)，\(AB = 2\)，\(BC = 2\sqrt{2}\)，\(SA = SB = \sqrt{3}\).

\begin{enumerate}
\item 证明：\(SA \perp BC\)；
\item 求直线 \(SD\) 与平面 \(SAB\) 所成角的大小.
\end{enumerate}

\bitmapfigure[max width=0.38\linewidth]{img/2007/syllabus_paper_1_liberal/q19_fig1.png}
\end{problem}

\begin{solution}
\begin{enumerate}
\item 建立空间直角坐标系. 取
\(B=(0,0,0)\)，\(C=(2\sqrt{2},0,0)\)，\(A=(\sqrt{2},\sqrt{2},0)\)，则
\(D=(3\sqrt{2},\sqrt{2},0)\). 因为平面 \(SBC\) 垂直于底面且包含直线 \(BC\)，可令
\(S=(u,0,v)\). 由 \(SA=SB=\sqrt{3}\)，有
\[
\begin{cases}
SB^2=u^2+v^2=3\\
SA^2=(u-\sqrt{2})^2+2+v^2=3
\end{cases}
\]
两式相减得 \(u=\sqrt{2}\)，再由 \(u^2+v^2=3\) 得 \(v^2=1\). 取坐标系正方向使 \(v=1\)，于是 \(S=(\sqrt{2},0,1)\).

此时 \(\overrightarrow{SA}=(0,\sqrt{2},-1)\)，\(\overrightarrow{BC}=(2\sqrt{2},0,0)\)，故
\(\overrightarrow{SA}\cdot\overrightarrow{BC}=0\)，从而 \(SA\perp BC\).
\item 由上面的坐标，\(\overrightarrow{SB}=(-\sqrt{2},0,-1)\)，\(\overrightarrow{SA}=(0,\sqrt{2},-1)\). 取平面 \(SAB\) 的一个法向量
\(\bs{n}=\overrightarrow{SA}\times\overrightarrow{SB}=(-\sqrt{2},\sqrt{2},2)\). 又
\(\overrightarrow{SD}=(2\sqrt{2},\sqrt{2},-1)\)，所以
\(\left|\overrightarrow{SD}\cdot\bs{n}\right|=4\)，\(\left|\overrightarrow{SD}\right|=\sqrt{11}\)，\(\left|\bs{n}\right|=2\sqrt{2}\).
设直线 \(SD\) 与平面 \(SAB\) 所成的角为 \(\theta\)，则
\(\sin\theta=\frac{\left|\overrightarrow{SD}\cdot\bs{n}\right|}{\left|\overrightarrow{SD}\right|\left|\bs{n}\right|}=\frac{\sqrt{22}}{11}\). 因此
\(\theta=\arcsin\frac{\sqrt{22}}{11}\).
\end{enumerate}
\end{solution}

\begin{problem}
设函数 \(f(x) = 2x^{3} + 3ax^{2} + 3bx + 8c\) 在 \(x = 1\) 及 \(x = 2\) 时取得极值.

\begin{enumerate}
\item 求 \(a\)，\(b\) 的值；
\item 若对于任意的 \(x \in [0,3]\)，都有 \(f(x) < c^{2}\) 成立，求 \(c\) 的取值范围.
\end{enumerate}
\end{problem}

\begin{solution}
\begin{enumerate}
\item \(f'(x)=6x^2+6ax+3b\). 因为函数在 \(x=1\)，\(x=2\) 时取得极值，所以
\(f'(1)=f'(2)=0\)，即
\[
\begin{cases}
2+2a+b=0\\
8+4a+b=0
\end{cases}
\]
两式相减得 \(a=-3\)，代回得 \(b=4\).
\item 代入上面的 \(a\)，\(b\)，得
\(f(x)=2x^3-9x^2+12x+8c\)，且
\(f'(x)=6(x-1)(x-2)\). 因此在 \([0,3]\) 上，函数分别在 \([0,1]\)，\([2,3]\) 上递增，在 \([1,2]\) 上递减. 计算得
\(f(0)=8c\)，\(f(1)=5+8c\)，\(f(2)=4+8c\)，\(f(3)=9+8c\).
所以 \([0,3]\) 上的最大值为 \(9+8c\). 题设要求对所有 \(x\in[0,3]\) 都有严格不等式，故
\(9+8c<c^2\)，即
\((c-9)(c+1)>0\). 因此 \(c<-1\) 或 \(c>9\).
\end{enumerate}
\end{solution}

\begin{problem}
设 \(\{a_{n}\}\) 是等差数列，\(\{b_{n}\}\) 是各项都为正数的等比数列，且 \(a_{1} = b_{1} = 1\)，\(a_{3} + b_{5} = 21\)，\(a_{5} + b_{3} = 13\).

\begin{enumerate}
\item 求 \(\{a_{n}\}\)，\(\{b_{n}\}\) 的通项公式；
\item 求数列 \(\left\{\frac{a_{n}}{b_{n}}\right\}\) 的前 \(n\) 项和 \(S_{n}\).
\end{enumerate}
\end{problem}

\begin{solution}
\begin{enumerate}
\item 设等差数列 \(\{a_n\}\) 的公差为 \(d\)，等比数列 \(\{b_n\}\) 的公比为 \(q\). 由首项条件
\(a_n=1+(n-1)d\)，\(b_n=q^{n-1}\)，且因各项为正数，\(q>0\).

由题设得
\[
\begin{cases}
2d+q^4=20\\
4d+q^2=12
\end{cases}
\]
令 \(t=q^2>0\)，则 \(2d+t^2=20\)，\(4d+t=12\). 消去 \(d\)，得
\(2t^2-t-28=0\)，即 \((t-4)(2t+7)=0\). 由于 \(t>0\)，所以 \(t=4\)，从而 \(q=2\)，\(d=2\). 因此
\(a_n=2n-1\)，\(b_n=2^{n-1}\).
\item 由 \(\frac{a_k}{b_k}=\frac{2k-1}{2^{k-1}}\)，有
\[
S_n=\sum_{k=1}^n\frac{2k-1}{2^{k-1}}=2\sum_{k=1}^n\frac{k}{2^{k-1}}-\sum_{k=1}^n\frac{1}{2^{k-1}}
\]
先由有限等比数列求和公式得
\[
\sum_{k=1}^n\frac{1}{2^{k-1}}
=\frac{1-\left(\frac{1}{2}\right)^n}{1-\frac{1}{2}}
=2-\frac{1}{2^{n-1}}
\]
再由
\(\sum_{k=1}^n r^k=\frac{r(1-r^n)}{1-r}\)，\((r\ne1)\)
两边对 \(r\) 求导，得
\[
\sum_{k=1}^nkr^{k-1}
=\frac{1-(n+1)r^n+nr^{n+1}}{(1-r)^2}
\]
令 \(r=\frac{1}{2}\)，可得
\[
\sum_{k=1}^n\frac{k}{2^{k-1}}
=\frac{1-\frac{n+1}{2^n}+\frac{n}{2^{n+1}}}{\left(1-\frac{1}{2}\right)^2}
=4-\frac{n+2}{2^{n-1}}
\]
代入 \(S_n\) 的表达式，得
\[
S_n=2\left(4-\frac{n+2}{2^{n-1}}\right)-\left(2-\frac{1}{2^{n-1}}\right)
=6-\frac{2n+3}{2^{n-1}}
\]
\end{enumerate}
\end{solution}

\begin{problem}
已知椭圆 \(\frac{x^2}{3} + \frac{y^2}{2} = 1\) 的左，右焦点分别为 \(F_1\)，\(F_2\). 过 \(F_1\) 的直线交椭圆于 \(B\)，\(D\) 两点，过 \(F_2\) 的直线交椭圆于 \(A\)，\(C\) 两点，且 \(AC \perp BD\)，垂足为 \(P\).

\begin{enumerate}
\item 设 \(P\) 点的坐标为 \((x_0, y_0)\). 证明：\(\frac{x_0^{2}}{3} + \frac{y_0^{2}}{2} < 1\)；
\item 求四边形 \(ABCD\) 的面积的最小值.
\end{enumerate}
\end{problem}

\begin{solution}
\begin{enumerate}
\item 椭圆的焦半距满足 \(c^2=3-2=1\)，所以
\(F_1=(-1,0)\)，\(F_2=(1,0)\). 因为 \(P\) 在直线 \(BD\) 和 \(AC\) 上，且 \(BD\perp AC\)，当 \(P\) 不与两焦点重合时，\(\overrightarrow{PF_1}\) 和 \(\overrightarrow{PF_2}\) 分别沿 \(BD\) 和 \(AC\)，所以
\[
\overrightarrow{PF_1}\cdot\overrightarrow{PF_2}=(-1-x_0)(1-x_0)+(-y_0)^2=x_0^2+y_0^2-1=0
\]
若 \(P=F_1\) 或 \(P=F_2\)，也有 \(x_0^2+y_0^2=1\). 因此无论哪种情况都有
\(x_0^2+y_0^2=1\)，于是
\(\frac{x_0^2}{3}+\frac{y_0^2}{2}=\frac{1}{2}-\frac{x_0^2}{6}\leqslant\frac{1}{2}<1\).
\item 先设直线 \(BD\) 不垂直于 \(x\) 轴，且斜率为 \(m\ne0\)，令 \(u=m^2>0\). 因为 \(BD\) 过 \(F_1=(-1,0)\)，其方程为 \(y=m(x+1)\). 与椭圆联立得
\[
(2+3u)x^2+6ux+3u-6=0
\]
该方程两根之差的绝对值为
\[
\frac{\sqrt{(6u)^2-4(2+3u)(3u-6)}}{2+3u}=\frac{4\sqrt{3(u+1)}}{2+3u}
\]
设两根为 \(x_1\)，\(x_2\). 由于直线 \(BD\) 的斜率为 \(m\)，两交点的纵坐标之差为 \(m(x_1-x_2)\)，所以
\[
BD=\sqrt{(x_1-x_2)^2+\bigl(m(x_1-x_2)\bigr)^2}
=\left|x_1-x_2\right|\sqrt{1+m^2}
=\frac{4\sqrt{3}(u+1)}{2+3u}
\]
直线 \(AC\) 的斜率为 \(-\frac{1}{m}\)，其斜率平方为 \(\frac{1}{u}\)，代入同一弦长公式，得
\[
AC=\frac{4\sqrt{3}\left(1+\frac{1}{u}\right)}{2+\frac{3}{u}}
=\frac{4\sqrt{3}(u+1)}{2u+3}
\]

四边形的两条对角线 \(AC\)，\(BD\) 垂直，故
\[
[ABCD]=\frac{1}{2}AC\cdot BD=\frac{24(u+1)^2}{(3u+2)(2u+3)}
\]
又
\[
\frac{(u+1)^2}{(3u+2)(2u+3)}-\frac{4}{25}
=\frac{(u-1)^2}{25(3u+2)(2u+3)}\geqslant0
\]
所以 \([ABCD]\geqslant\frac{96}{25}\). 当 \(u=1\)，即 \(m=\pm1\) 时取等号.

再看未包含在上述斜率讨论中的特殊方向：若 \(BD\) 水平，则 \(BD\) 是椭圆的长轴，故 \(BD=2\sqrt{3}\)；直线 \(AC\) 为过 \(F_2\) 的竖直弦，令 \(x=1\) 代入椭圆方程，得 \(y=\pm\frac{2}{\sqrt{3}}\)，所以 \(AC=\frac{4}{\sqrt{3}}\)，面积为 \(4\). 若 \(BD\) 竖直，则令 \(x=-1\) 得 \(BD=\frac{4}{\sqrt{3}}\)，而 \(AC\) 为长轴，故 \(AC=2\sqrt{3}\)，面积仍为 \(4\). 因为 \(4>\frac{96}{25}\)，所以四边形 \(ABCD\) 的面积最小值为 \(\frac{96}{25}\).
\end{enumerate}
\end{solution}
